\section{Anlageurteil}
\label{sec:urteil}

\subsection{Kurzfristiges Potenzial, null bis vierundzwanzig Monate}

% FAKTENBASIS 6.1
%  - Organisches Wachstum H1 2026 \HJOrganisch %, Q2 allein \QzweiOrganisch %
%  - Berichteter Umsatz H1 2026 +\HJUmsatzDelta % (Konsolidierungseffekt)
%  - Kurs \KursSchlussZFVI USD, \KursAbstandHoch % unter dem Hoch 2025
%  - Aktienrueckkauf H1 2026 \HJEigeneAktienZuwachs Mio. USD
%  - Verschuldungsgrad 2026 \FlowGradStart, Zinsdeckung \Zinsdeckung
\bk{Die kommenden vier bis sechs Quartale entscheiden sich am organischen
Wachstum und an nichts sonst. Der berichtete Umsatz wird bis zum dritten
Quartal 2026 weiter zweistellig wachsen, weil die \"Ubernahme sich erst ab
August 2026 im Vorjahresvergleich neutralisiert; dieses Wachstum ist ein
Konsolidierungseffekt und tr\"agt kein Urteil. Der Fr\"uhindikator daneben
ist im ersten Halbjahr 2026 mit \Pct{\HJOrganisch} negativ geworden, im
zweiten Quartal allein mit \Pct{\QzweiOrganisch}. Solange er das bleibt,
gibt es keinen Grund, eine Erholung des Kurses zu erwarten, und der
R\"uckkauf von \MioUSD{\HJEigeneAktienZuwachs} im ersten Halbjahr ist
allenfalls eine St\"utze, kein Wendepunkt.

Best\"atigen oder widerlegen l\"asst sich das an einer einzigen, quartalsweise
ver\"offentlichten Zahl. Kehrt die organische Rate in den positiven Bereich
zur\"uck, ist der Einbruch als Verdauung der Integration lesbar. Bleibt sie
zwei weitere Quartale negativ, ist sie ein Zeichen f\"ur Kundenabgang aus
der \"Ubernahme selbst; das w\"are der teuerste denkbare Verlauf, weil
\Faktor{\KaufpreisJeUmsatz}-facher Jahresumsatz genau f\"ur diese Kunden
bezahlt wurde.}

\subsection{Langfristiges Potenzial, drei bis sieben Jahre}

% FAKTENBASIS 6.2
%  - Verschuldungsgrad faellt \FlowGradStart -> \FlowGradEnde (mittleres
%    Szenario), im schlechtesten Szenario auf \FlowGradSzenarioA
%  - Sachinvestitionen \FlowSachinvest Mio. USD gegen EBITDAC \FlowEbitdacEnde
%  - IRR 15 Jahre \FlowIRRBFuenfzehn % gegen Massstab \EKRenditeDurchschnitt %
%  - Goodwill \Goodwill Mio. USD = \GoodwillQuote % der Bilanzsumme,
%    \GoodwillEKDeckung-faches Eigenkapital
\bk{Strukturell ist der Fall besser als der Kursverlauf nahelegt. Ein
Versicherungsmakler bindet kaum Kapital: Sachinvestitionen von
\MioUSD{\FlowSachinvest} stehen einem bereinigten EBITDAC von
\MioUSD{\FlowEbitdacEnde} gegen\"uber. Das Ergebnis wird deshalb fast
vollst\"andig zu Bargeld, und der Verschuldungsgrad f\"allt selbst im
ung\"unstigsten der drei Szenarien von \Faktor{\FlowGradStart} auf
\Faktor{\FlowGradSzenarioA}. Die \"Ubernahme ist finanzierbar; die Frage ist
nur, ob sie ihren Preis wert war.

Damit sie es ist, muss eine von zwei Bedingungen halten. Entweder kehrt das
organische Wachstum auf ein mittleres einstelliges Niveau zur\"uck, dann
liegt der interne Zinsfu{\ss} \"uber f\"unfzehn Jahre bei
\Pct{\FlowIRRBFuenfzehn} und damit ungef\"ahr auf dem Ma{\ss}stab der eigenen
Eigenkapitalrendite. Oder das Unternehmen setzt seinen freien Cash-Flow
erneut in Zuk\"aufe um und verdient daran mehr als die
\Faktor{\KaufpreisJeUmsatz}-fache Umsatzbewertung kostet. Das zweite ist
das tats\"achliche Gesch\"aftsmodell und in dieser Rechnung bewusst nicht
abgebildet.}

\subsection{Zentrale Abh\"angigkeit}

\bk{Der Fall ruht auf einer einzigen Gr\"o{\ss}e: dem organischen
Umsatzwachstum. Sie erkl\"art den Kursverlauf der vergangenen zwei Jahre
(Teil~\ref{sec:kurs}), sie tr\"agt die Spannweite des internen Zinsfu{\ss}es
von \Pct{\FlowIRRAFuenf} bis \Pct{\FlowIRRCFuenf} \"uber f\"unf Jahre
(Tabelle~\ref{tab:flowirr}), und sie ist zugleich die Gr\"o{\ss}e, an der
sich zeigen wird, ob der Kaufpreis der \"Ubernahme berechtigt war. Die
Verschuldung ist demgegen\"uber zweitrangig: Sie ist hoch, aber der
Cash-Flow tr\"agt sie in jedem gerechneten Szenario.

Eine zweite, stillere Abh\"angigkeit steht in der Bilanz. Der Firmenwert von
\MioUSD{\Goodwill} entspricht \Pct{\GoodwillQuote} der Bilanzsumme und dem
\Faktor{\GoodwillEKDeckung}-fachen des Eigenkapitals. Eine Wertberichtigung
darauf w\"urde das Eigenkapital unmittelbar treffen. Sie ist keine
Zahlungsgr\"o{\ss}e und ver\"andert weder Tilgungsf\"ahigkeit noch internen
Zinsfu{\ss}, w\"are aber das sichtbare Eingest\"andnis, dass zu teuer gekauft
wurde.}

\subsection{Was das Urteil \"andern w\"urde}

\begin{table}[htbp]
\centering
\caption{Beobachtbare Ausl\"oser, die das Urteil in die eine oder andere
Richtung verschieben.}
\label{tab:ausloeser}
\begin{tabularx}{\linewidth}{lX}
\toprule
Richtung & Ausl\"oser\\
\midrule
Zum Besseren & Organische Rate kehrt \"uber zwei aufeinanderfolgende Quartale in den positiven Bereich zur\"uck.\\
             & Bereinigte EBITDAC-Marge steigt \"uber \Pct{\EBITDACberMarge}, weil Integrationskosten auslaufen.\\
             & Verschuldungsgrad unterschreitet vor Ende 2027 den Wert, mit dem
               die \"Ubernahme startete (\Faktor{\VerschuldungsgradProForma}).\\
\midrule
Zum Schlechteren & Organische Rate bleibt zwei weitere Quartale negativ.\\
             & Wertberichtigung auf den Firmenwert aus dem Accession-Erwerb.\\
             & Erneuter gro{\ss}er Zukauf, bevor der Verschuldungsgrad unter das Vor\"ubernahmeniveau von \Faktor{\VerschuldungsgradVJ} zur\"uckgefallen ist.\\
\bottomrule
\end{tabularx}
\end{table}

\FloatBarrier

\subsection{Vertrauensgrad}

\bk{Der Vertrauensgrad dieser Analyse ist mittel, und die Grenze setzen die
in Abschnitt~\ref{sec:luecken} und unten benannten L\"ucken.} Belastbar sind
die historischen Zahlen: Jede Angabe der Teile~\ref{sec:organisch} bis
\ref{sec:kurs} ist auf die gedruckte Seite einer Prim\"arquelle
zur\"uckgef\"uhrt und maschinell gegen diese gepr\"uft. Weniger belastbar
ist die Fortschreibung, und zwar aus drei benannten Gr\"unden: der
unterstellten konstanten Marge, der ausgeschlossenen weiteren Zuk\"aufe und
des unver\"anderten Ausstiegsvielfachen. Die ersten beiden sind in
Teil~\ref{sec:flow} als Annahmen ausgeschrieben, das dritte ist in
Tabelle~\ref{tab:flowmultiple} variiert.

\subsection*{Nicht verf\"ugbare Angaben}
\addcontentsline{toc}{subsection}{Nicht verf\"ugbare Angaben}

Folgende Gr\"o{\ss}en waren aus den Prim\"arquellen nicht zu belegen und
werden deshalb nicht angegeben und nicht gesch\"atzt.

\begin{enumerate}
\item \textbf{Ratings und Kursziele.} Die Investor-Relations-Seite nennt
\num{\AnalystenHaeuser} abdeckende H\"auser mit Namen, ver\"offentlicht aber
weder Empfehlungen noch Kursziele.\quelleWeb{Brown \& Brown, Inc., Analyst Coverage}{quelle:analysten}{28.08.2026}\quelleMerken{s17webbrownbrowni} Ein
Konsensrating wird deshalb nicht gebildet.
\item \textbf{Unternehmensprognose.} Als US-amerikanischer Emittent
ver\"offentlicht das Unternehmen im Jahresbericht keine Umsatz- oder
Margenbandbreite; ein Abgleich von Ausblick und Ist ist nicht durchf\"uhrbar
(Abschnitt~\ref{sec:ausblick}).
\item \textbf{Kapitalkosten.} Beta, risikofreier Zins und
Marktrisikopr\"amie stehen in keiner der Prim\"arquellen. Statt eines
gesch\"atzten Kapitalkostensatzes dient die belegte Eigenkapitalrendite von
\Pct{\EKRenditeDurchschnitt} als Vergleichsma{\ss}stab.
\item \textbf{Weitere Gro{\ss}aktion\"are.} Das Proxy Statement weist neben
Vanguard und Capital World keinen Halter oberhalb von f\"unf Prozent aus.
Ob weitere institutionelle Investoren knapp darunter liegen, ist der
Prim\"arquelle nicht zu entnehmen.
\item \textbf{Getrennter Ergebnisbeitrag von Accession.} Der Jahresbericht
weist Umsatz und Ergebnis der Zuk\"aufe 2025 nur zusammengefasst
aus.\quelleGB{2025}{63}\quelleMerken{s18gb} Der in Abschnitt~\ref{sec:accession}
genannte Jahresumsatz von \MioUSD{\ZukaeufeUmsatzJahr} ist deshalb eine aus
zwei ver\"offentlichten Angaben hergeleitete Gr\"o{\ss}e f\"ur alle
\num{\AkquisitionenAnzahl} Zuk\"aufe zusammen und keine Angabe zu Accession
allein; die Formel steht dort.
\end{enumerate}
